\begin{pastexam}[label={px:ncku-114-cce-ds-12}]{成大電通所 114 學年度 單選題第 12 題}
（單選；每題 5 分，答錯倒扣 2 分）

Continue from the problem 11, instead of chaining, we adopt open addressing to handle overflow, where each bucket
has one slot. The resulting position of data is determined by h(key)= \{initial position + $i$* h$_2$(key) \}\%11,
where h$_2$(key) = key \% 7 and $i$ denotes the iteration of the overflow. Please show the position of 58 in the
table.\\
a) 7\quad b) 5\quad c) 10\quad d) 3\quad e) None of the above.
\end{pastexam}

\begin{pastsoln}
答案：a) 7。

這是雙重雜湊（資結書第~12.3~節）：第 $i$ 次探測的位置是 $(\mathit{key} \bmod 11 + i \cdot (\mathit{key} \bmod 7)) \bmod 11$，
$i = 0, 1, 2, \ldots$，第一個空格就放進去。 鍵依第 11 題的順序插入，58 之前的五個鍵：
\begin{center}
\begin{tabular}{lccl}
\toprule
key & key \% 11 & key \% 7 & 探測位置 \\
\midrule
3 & 3 & & 3 \\
41 & 8 & & 8 \\
15 & 4 & & 4 \\
36 & 3 & 1 & 3（被 3 佔）$\to$ 4（被 15 佔）$\to$ 5，放在 5 \\
74 & 8 & 4 & 8（被 41 佔）$\to$ $12 \bmod 11 = 1$，放在 1 \\
58 & 3 & 2 & 3（被 3 佔）$\to$ 5（被 36 佔）$\to$ 7，放在 7 \\
\bottomrule
\end{tabular}
\end{center}
58 放在位置 7，選 a)。

附帶一提，這個 $h_2$ 設計得不好：$h_2$ 可能是 0。 接下來的 91 就是例子：$91 \bmod 11 = 3$ 已被佔用，而 $91 \bmod 7 = 0$，
每次探測都停在 3，永遠找不到空格。 雙重雜湊的第二個函數不能是 0，而且要和表的大小互質（資結書第~12.3~節）；表的大小 11 是質數時，
常見的取法是讓 $h_2$ 的值落在 $1$ 到 $10$ 之間，例如 $h_2(k) = 1 + (k \bmod 7)$。 這不影響 58 的答案。
\end{pastsoln}
